\section{Identical scalar orders in the binary experiment}\label{sec:valuative}
The example in Section~\ref{sec:separation} keeps sharp marginal ranges
fixed. We now keep a different, much larger valuative object fixed inside
the full experiment of Section~\ref{sec:model}: all real analytic arc
orders of the observation discrepancy and every polynomial scalar
function of the target coefficients. Leading units are not included in
this datum. The distinction is necessary, since
Theorem~\ref{thm:reconstruction} shows how richer scalar values recover
the joint set.

Take degree one, $D=1$, $\pi=1/10$, clocks $(2,3,4)$ and weights $1/3$.
Consider the two strictly interior full-rank centres $\theta_e$, with
$e\in\{1/4,1/8\}$, given by
\begin{equation}\label{eq:binary-two-centres}
 f=z-\tfrac14,\quad g=z-\tfrac34,\quad \alpha=\tfrac12,
 \quad U_e=\begin{pmatrix}\frac12+e&\frac12-e\\
                          \frac12-e&\frac12+e\end{pmatrix},
 \quad V=\begin{pmatrix}\frac14&\frac34\\\frac34&\frac14\end{pmatrix}.
\end{equation}
All competitors still range over the original closed binary model; no
parameter is held known. In a common small open coordinate ball write
\[
 z=(\dot r,\dot s,\dot\alpha,\dot u_1,\dot u_2,
                     \dot v_1,\dot v_2).
\]
The first entries of the binary channel columns are the independent
channel coordinates. Let $G_e(z)$ be squared Hellinger discrepancy from
$\theta_e$, and let $\Delta c(z)$ be the two target coefficient
increments. The latter polynomial map is the same at both centres.

\begin{theorem}[Valuative equivalence without metric equivalence]\label{thm:binary-separation}
For every nonconstant real analytic arc $z(\tau)$ through zero,
\begin{equation}\label{eq:all-arc-orders}
 \operatorname{ord}_\tau G_{1/4}(z(\tau))
 =\operatorname{ord}_\tau G_{1/8}(z(\tau))
 =2\min_i\operatorname{ord}_\tau z_i(\tau).
\end{equation}
For every polynomial $H$ in the target coefficient increments, including
all mixed polynomial combinations, the orders of
$H(\Delta c(z(\tau)))$ agree at the two centres, with the same convention
for identically zero functions. The corresponding accessible real
divisorial order data can be computed on one common real monomial
presentation and agree as well.

Nevertheless the two whole-model leading spectral sets are distinct.
Their exponents are one, and their squared leading diameters are
\begin{equation}\label{eq:binary-diameters}
 C_{1/4}^{\,2}=\frac{129609775953}{25280},\qquad
 C_{1/8}^{\,2}=\frac{10058351786055}{1396384}.
\end{equation}
Thus the complete scalar arc-order information just specified does not
determine the joint metric initial fibre, even within the principal
binary observation model.
\end{theorem}
\begin{proof}
Both channels are invertible and the component polynomials are coprime.
Lemma~\ref{lem:fullrank} implies that the observation derivative in
these seven coordinates is injective. Since all probability cells are
positive, the square-root embedding is analytic and
\[
 G_e(z)=\tfrac14 z^{\mathsf T}H_ez+O(\|z\|^3),\qquad H_e>0.
\]
It follows either by this expansion or by two-sided comparison with
$\|z\|^2$ that \eqref{eq:all-arc-orders} holds. The coefficient functions
are literally identical in the common coordinates, proving their order
assertion for every $H$.

For the divisorial assertion, first blow up the maximal ideal at zero.
On every real affine chart the pullback of $G_e$ is the square of an
exceptional coordinate times a strictly positive analytic unit near the
real exceptional locus. Indeed the leading factor is the positive
quadratic form $H_e$ evaluated on a nonzero real projective direction.
These units remain nonzero under every further real modification.
Principalize any finite specified collection of the common coefficient
functions on this same space. The exceptional orders and their real
accessibility are the same for both experiments. Alternatively every
accessible divisor has a real transversal, so its observation order
already follows from \eqref{eq:all-arc-orders}. The assertion concerns
real accessible data; no equality of complex zero divisors of the two
positive quadratic forms is used. Every additional finite collection of
polynomial combinations is handled in the same way.

We compute the metric sets. Let $S_e$ be the $12$-by-$7$ score matrix
from \eqref{eq:score}, at the three clocks, in the coordinate order above.
Set
\[
 H_e=S_e^{\mathsf T}\diag\bigl((3P_{e,je})^{-1}\bigr)S_e,
 \qquad M_e=(H_e^{-1})_{\{1,2\},\{1,2\}}.
\]
Profiling the five free nuisance coordinates in the positive definite
quadratic form shows that the leading pair of root increments is exactly
\begin{equation}\label{eq:binary-ellipse}
 L_e=\{x\in\R^2:x^{\mathsf T}M_e^{-1}x\le4\}.
\end{equation}
For completeness, the lower bound on $G_e$ bounds all competitor
coordinates by $O(t)$. Division by $t$ and Taylor expansion give the
outer inclusion in \eqref{eq:binary-ellipse}. Every vector in its
interior is realized by its linear parameter arc, which is feasible
because the centre is interior. Boundary vectors are limits of their
radial contractions. This proves Hausdorff convergence. Exact uniqueness
and compactness, from Lemma~\ref{lem:fullrank}, exclude competitors
outside this coordinate neighbourhood at sufficiently small radii.
The limit is consequently a whole-model limit.

Here the entries of the matrix are obtained by differentiating the
explicit rational functions
\[
 P_{j,ab}=
 \frac{\alpha(T_j-r)u_{1a}v_{1b}
 +(1-\alpha)(T_j-s)u_{2a}v_{2b}}
 {\alpha(T_j-r)+(1-\alpha)(T_j-s)}.
\]
Rational elimination gives
\begin{align}
 M_{1/4}&=\frac1{404480}
 \begin{pmatrix}129609775953&79361335413\\
                 79361335413&48691215153\end{pmatrix},\label{eq:M-quarter}\\
 M_{1/8}&=\frac1{22342144}
 \begin{pmatrix}10058351786055&6157532670525\\
                 6157532670525&3781629625095\end{pmatrix}.
                                                        \label{eq:M-eighth}
\end{align}
These equalities can be checked solely by forming $H_e$ from the preceding
rational display and multiplying its inverse; no fitted numerical
covariance is involved. For separated real roots the bottleneck metric
is the maximum norm of their ordered increments. A centrally symmetric
ellipsoid \eqref{eq:binary-ellipse} has maximum-norm diameter
$4\max_i\sqrt{(M_e)_{ii}}$. The first diagonal entry is the larger
in each displayed matrix. This proves \eqref{eq:binary-diameters}.
The values are unequal, so neither the leading root sets nor their
coefficient images can be equal.
\end{proof}

The example fixes all orders of the target scalar functions, not their
sharp leading values. Those values contain the positive units erased
by valuation. It therefore complements the sharp-marginal example
without claiming that arbitrary scalar information is insufficient.
The relative real atlas remains useful for finding the common weights;
its units and its joint relation are necessary for recovering the metric
section at those weights.
